\subsection{Отображение и взаимодействие}

После проецирования координаты переводятся в координаты области рисования: начало координат размещается в центре, вертикальное направление согласуется с экранной системой, а к изображению применяется общий масштаб. Масштаб определяется меньшей стороной области рисования, поэтому при изменении размеров окна пропорции фигуры сохраняются. Увеличенное изображение ограничивается границами этой области.

В режиме заливки отображаются только грани, обращённые к камере. Для определения видимости учитываются направление нормали грани и направление на камеру. Для ортографической проекции направление наблюдения одинаково для всех граней, а для перспективной определяется относительно центра конкретной грани. Это позволяет корректно менять видимые поверхности при вращении.

Видимые грани сортируются от дальних к ближним и передаются для рисования как выпуклые двумерные полигоны. Для текущего одиночного выпуклого объекта отсечение обратных граней достаточно для устранения скрытых поверхностей. Буфер глубины для полигонов модели не используется. Применение этого подхода к нескольким пересекающимся объектам потребовало бы другого решения.

Треугольные грани окрашены в оранжевый цвет, а шестиугольные — в синий. Яркость зависит от ориентации грани относительно фиксированного источника света. Контуры выделены тёмными линиями. В каркасном режиме выводятся все 18 рёбер, включая скрытые, что позволяет рассмотреть внутреннюю структуру проекции.

Поворот по трём осям регулируется отдельными ползунками. Перетаскивание в области рисования изменяет поворот вокруг осей X и Y. Колёсико мыши и ползунок Zoom меняют масштаб отображения. Кнопка Reset view восстанавливает исходные углы и масштаб. Флажки переключают проекцию, каркас и автоматическое вращение. Скорость автоматического вращения учитывает время между кадрами.
